% ============================================================================
% CHAPTER 6: TROUBLESHOOTING
% ============================================================================

\chapter{Troubleshooting Guide}

\section{Common Issues \& Solutions}

\subsection{Services Won't Start}

\textbf{Problem:} Docker services fail to start or show "unhealthy" status.

\textbf{Diagnosis:}
\begin{commandbox}
\begin{lstlisting}[style=bashstyle]
# Check service status
docker-compose ps

# View specific service logs
docker-compose logs <service-name>

# Check resource usage
docker stats
\end{lstlisting}
\end{commandbox}

\textbf{Common Causes \& Solutions:}

\begin{enumerate}
    \item \textbf{Port Already in Use}
    \begin{lstlisting}[style=bashstyle]
# Find process using port 3000
lsof -i :3000  # macOS/Linux
netstat -ano | findstr :3000  # Windows

# Kill process or change port in docker-compose.yml
\end{lstlisting}

    \item \textbf{Insufficient Memory}
    \begin{itemize}
        \item Increase Docker Desktop memory limit (Settings → Resources)
        \item Recommended: 8 GB minimum, 16 GB ideal
    \end{itemize}

    \item \textbf{Network Conflicts}
    \begin{lstlisting}[style=bashstyle]
# Remove old networks
docker network prune

# Recreate with clean network
docker-compose down
docker-compose up -d
\end{lstlisting}

    \item \textbf{Database Not Ready}
    \begin{lstlisting}[style=bashstyle]
# Wait for database health check
docker-compose up -d postgres elasticsearch neo4j
sleep 30  # Wait for initialization

# Then start application services
docker-compose up -d api-gateway frontend
\end{lstlisting}
\end{enumerate}

\subsection{Classification Not Working}

\textbf{Problem:} Documents uploaded but not classified, or classification fails.

\textbf{Diagnosis Steps:}
\begin{commandbox}
\begin{lstlisting}[style=bashstyle]
# 1. Check PV Classification Service logs
docker-compose logs pv-classification-service | tail -50

# 2. Verify service is accessible
curl http://localhost:8012/health

# 3. Test classification endpoint directly
curl -X POST http://localhost:8012/api/v1/classify \
  -H "Content-Type: application/json" \
  -d '{"title":"Test","abstract":"Patient developed rash"}'
\end{lstlisting}
\end{commandbox}

\textbf{Common Causes:}

\begin{enumerate}
    \item \textbf{Missing OpenAI API Key}
    \begin{itemize}
        \item LLM tier will fail without API key
        \item Solution: Set \texttt{OPENAI\_API\_KEY} in \texttt{.env}
        \item Alternative: Use Ollama local LLM (no API key needed)
    \end{itemize}

    \item \textbf{BERT Model Not Loaded}
    \begin{lstlisting}[style=bashstyle]
# Check if model is downloading
docker-compose logs pv-classification-service | grep "Downloading"

# Model downloads on first run (~1.5 GB)
# Wait 5-10 minutes for initial download
\end{lstlisting}

    \item \textbf{Ensemble Configuration Error}
    \begin{lstlisting}[style=pythonstyle]
# Verify ensemble initialization in logs
docker-compose logs pv-classification-service | grep "Ensemble"

# Expected:
# "Ensemble classifier initialized successfully"
# "Models loaded: ['rule_based', 'bert', 'llm']"
\end{lstlisting}

    \item \textbf{Port Mapping Mismatch}
    \begin{itemize}
        \item Verify \texttt{docker-compose.yml} port mapping: \texttt{"8012:8002"}
        \item Internal service runs on port 8002
        \item External access on port 8012
    \end{itemize}
\end{enumerate}

\subsection{Knowledge Graph Not Displaying}

\textbf{Problem:} Knowledge Hub shows empty graph or fails to load.

\textbf{Diagnosis:}
\begin{commandbox}
\begin{lstlisting}[style=bashstyle]
# 1. Check Neo4j is running
curl http://localhost:7474

# 2. Verify node count
docker-compose exec neo4j cypher-shell -u neo4j -p password \
  "MATCH (n) RETURN count(n) as NodeCount;"

# 3. Check API Gateway logs
docker-compose logs api-gateway | grep "knowledge-hub"

# 4. Test API endpoint
curl -H "Authorization: Bearer <token>" \
  http://localhost:8000/api/v1/knowledge-hub/graph
\end{lstlisting}
\end{commandbox}

\textbf{Common Causes:}

\begin{enumerate}
    \item \textbf{No Documents Processed}
    \begin{itemize}
        \item Upload and classify at least 10 documents first
        \item Entity extraction must be enabled
    \end{itemize}

    \item \textbf{Neo4j Not Populated}
    \begin{lstlisting}[style=bashstyle]
# Check if entities were synced to Neo4j
docker-compose logs api-gateway | grep "Neo4j sync"

# Manually trigger sync (if available)
curl -X POST http://localhost:8000/api/v1/sync-neo4j
\end{lstlisting}

    \item \textbf{Browser Cache Issue}
    \begin{itemize}
        \item Force refresh: Ctrl+Shift+R (Windows/Linux) or Cmd+Shift+R (macOS)
        \item Clear browser cache completely
        \item Try incognito/private mode
    \end{itemize}

    \item \textbf{WebGL Not Supported}
    \begin{itemize}
        \item Check browser WebGL support: \texttt{https://get.webgl.org/}
        \item Update graphics drivers
        \item Switch to 2D Canvas mode in UI
    \end{itemize}
\end{enumerate}

\subsection{Frontend Not Loading}

\textbf{Problem:} Blank page, "502 Bad Gateway", or React errors.

\textbf{Diagnosis:}
\begin{commandbox}
\begin{lstlisting}[style=bashstyle]
# 1. Check frontend service
docker-compose logs frontend

# 2. Check nginx errors
docker-compose exec frontend cat /var/log/nginx/error.log

# 3. Verify build completed successfully
docker-compose logs frontend | grep "Compiled successfully"

# 4. Test direct access to API Gateway
curl http://localhost:8000/health
\end{lstlisting}
\end{commandbox}

\textbf{Solutions:}

\begin{enumerate}
    \item \textbf{Rebuild Frontend}
    \begin{lstlisting}[style=bashstyle]
# Rebuild with no cache
docker-compose build --no-cache frontend
docker-compose up -d frontend

# Force browser refresh
# Press Ctrl+Shift+R (hard reload)
\end{lstlisting}

    \item \textbf{Fix API Gateway Connection}
    \begin{itemize}
        \item Check \texttt{REACT\_APP\_API\_URL} in \texttt{.env}
        \item Should be: \texttt{http://localhost:8000}
        \item Restart frontend after changing environment variables
    \end{itemize}

    \item \textbf{Clear Browser Cache}
    \begin{lstlisting}[style=bashstyle]
# Chrome DevTools → Application → Clear Storage → Clear site data
# Or use Incognito mode to test
\end{lstlisting}

    \item \textbf{Check Console Errors}
    \begin{itemize}
        \item Open browser DevTools (F12)
        \item Check Console tab for JavaScript errors
        \item Check Network tab for failed API calls
    \end{itemize}
\end{enumerate}

\subsection{High Memory Usage}

\textbf{Problem:} System becomes slow, services crash due to OOM (Out of Memory).

\textbf{Diagnosis:}
\begin{commandbox}
\begin{lstlisting}[style=bashstyle]
# Check container memory usage
docker stats --no-stream

# Expected output:
# NAME                 MEM USAGE / LIMIT     MEM %
# elasticsearch        3.8GB / 4GB           95%   <-- High!
# postgres             450MB / 2GB           22%
# pv-classification    2.1GB / 4GB           52%
\end{lstlisting}
\end{commandbox}

\textbf{Solutions:}

\begin{enumerate}
    \item \textbf{Increase Docker Memory Limit}
    \begin{itemize}
        \item Docker Desktop → Settings → Resources
        \item Increase memory to 16 GB minimum
        \item Restart Docker Desktop
    \end{itemize}

    \item \textbf{Optimize Elasticsearch}
    \begin{lstlisting}[style=bashstyle]
# Reduce heap size if limited memory
# Edit docker-compose.yml:
environment:
  - "ES_JAVA_OPTS=-Xms1g -Xmx1g"  # Reduce from 2g

# Restart Elasticsearch
docker-compose restart elasticsearch
\end{lstlisting}

    \item \textbf{Clear Unused Data}
    \begin{lstlisting}[style=bashstyle]
# Remove old indices
curl -X DELETE "http://localhost:9200/old_index_name"

# Vacuum PostgreSQL
docker-compose exec postgres psql -U postgres -d document_intelligence -c "VACUUM FULL;"

# Clean Docker system
docker system prune -a --volumes
\end{lstlisting}

    \item \textbf{Scale Down Services}
    \begin{itemize}
        \item Stop non-essential services during peak processing
        \item Example: Stop Grafana/Prometheus during batch ingestion
    \end{itemize}
\end{enumerate}

\subsection{Slow Classification Performance}

\textbf{Problem:} Classification takes too long (>10 seconds per document).

\textbf{Expected Performance:}
\begin{itemize}
    \item Rule-based only: 50-100ms
    \item Rule + BERT: 800-1200ms
    \item Full ensemble (with LLM): 2-4 seconds
\end{itemize}

\textbf{Optimization Steps:}

\begin{enumerate}
    \item \textbf{Use GPU Acceleration (BERT)}
    \begin{lstlisting}[style=pythonstyle]
# docker-compose.yml - Add GPU support
services:
  pv-classification-service:
    deploy:
      resources:
        reservations:
          devices:
            - driver: nvidia
              count: 1
              capabilities: [gpu]
\end{lstlisting}

    \item \textbf{Increase Confidence Thresholds}
    \begin{itemize}
        \item Higher thresholds = fewer LLM invocations
        \item Edit: \texttt{services/pv-classification-service/config.py}
        \item Increase \texttt{RULE\_CONFIDENCE\_THRESHOLD} from 0.85 to 0.90
    \end{itemize}

    \item \textbf{Use Ollama Instead of OpenAI}
    \begin{itemize}
        \item Local LLM = no API latency
        \item Requires more CPU/Memory but faster
        \item Install Ollama: \texttt{curl https://ollama.ai/install.sh | sh}
    \end{itemize}

    \item \textbf{Batch Processing}
    \begin{lstlisting}[style=pythonstyle]
# Use batch classification endpoint
POST /api/v1/pv-classification/classify/batch
{
  "documents": [doc1, doc2, ..., doc100]
}
# Processes 100 documents in parallel
\end{lstlisting}

    \item \textbf{Enable Caching}
    \begin{itemize}
        \item Classification results cached in Redis (30 min)
        \item Duplicate documents return cached result
        \item No re-classification needed
    \end{itemize}
\end{enumerate}

\section{Error Messages \& Meanings}

\subsection{Backend Errors}

\begin{table}[h]
\centering
\caption{Common Backend Error Messages}
\begin{longtable}{p{0.35\textwidth}p{0.6\textwidth}}
\toprule
\textbf{Error Message} & \textbf{Cause \& Solution} \\
\midrule
\texttt{Elasticsearch connection failed} & 
Elasticsearch not running or wrong URL. Check \texttt{ELASTICSEARCH\_URL} in \texttt{.env}. \\
\addlinespace
\texttt{Neo4j authentication failed} & 
Wrong password in \texttt{NEO4J\_AUTH}. Verify credentials match. \\
\addlinespace
\texttt{OpenAI API key not found} & 
Missing \texttt{OPENAI\_API\_KEY} in \texttt{.env}. Set key or use Ollama. \\
\addlinespace
\texttt{Model not found: bert} & 
BERT model not downloaded. Wait for download or check disk space. \\
\addlinespace
\texttt{Classification timeout} & 
Document too large or LLM slow. Increase timeout or split document. \\
\addlinespace
\texttt{JWT token expired} & 
Access token expired (30 min). Refresh token or re-login. \\
\addlinespace
\texttt{Insufficient permissions} & 
User role lacks permission. Check RBAC settings. \\
\addlinespace
\texttt{Database connection pool exhausted} & 
Too many concurrent connections. Increase pool size or scale database. \\
\bottomrule
\end{longtable}
\end{table}

\subsection{Frontend Errors}

\begin{table}[h]
\centering
\caption{Common Frontend Error Messages}
\begin{longtable}{p{0.35\textwidth}p{0.6\textwidth}}
\toprule
\textbf{Error Message} & \textbf{Cause \& Solution} \\
\midrule
\texttt{Network Error: Failed to fetch} & 
API Gateway not reachable. Check if backend is running. \\
\addlinespace
\texttt{Unexpected token '<' in JSON} & 
Receiving HTML instead of JSON. Check API endpoint URL. \\
\addlinespace
\texttt{WebGL not supported} & 
Browser/GPU doesn't support WebGL. Update drivers or use 2D mode. \\
\addlinespace
\texttt{Cannot read property 'map' of undefined} & 
API returned unexpected data structure. Check API response format. \\
\addlinespace
\texttt{401 Unauthorized} & 
Token missing or invalid. Re-login to get fresh token. \\
\addlinespace
\texttt{CORS policy blocked} & 
API Gateway CORS not configured. Add frontend origin to allowed list. \\
\bottomrule
\end{longtable}
\end{table}

\section{Debugging Techniques}

\subsection{Enable Debug Logging}

\textbf{API Gateway:}
\begin{lstlisting}[style=pythonstyle]
# Edit services/api-gateway/main.py
import logging
logging.basicConfig(level=logging.DEBUG)

# Or set environment variable
LOG_LEVEL=DEBUG docker-compose up -d api-gateway
\end{lstlisting}

\textbf{PV Classification Service:}
\begin{lstlisting}[style=pythonstyle]
# Edit services/pv-classification-service/main.py
import logging
logging.getLogger("uvicorn").setLevel(logging.DEBUG)
logging.getLogger("pv_classification").setLevel(logging.DEBUG)
\end{lstlisting}

\subsection{Interactive Debugging}

\textbf{Attach to Running Container:}
\begin{commandbox}
\begin{lstlisting}[style=bashstyle]
# Open shell in container
docker-compose exec api-gateway /bin/bash

# Install debugging tools (temporary)
apt-get update && apt-get install -y curl jq

# Test internal service connections
curl http://elasticsearch:9200/_cluster/health
curl http://pv-classification-service:8002/health
\end{lstlisting}
\end{commandbox}

\textbf{Python Debugger (pdb):}
\begin{lstlisting}[style=pythonstyle]
# Add breakpoint in code
import pdb; pdb.set_trace()

# Attach to container with stdin
docker-compose run --service-ports api-gateway python main.py
\end{lstlisting}

\subsection{Network Diagnostics}

\textbf{Test Service Connectivity:}
\begin{commandbox}
\begin{lstlisting}[style=bashstyle]
# Check if service is listening
docker-compose exec api-gateway netstat -tuln | grep 8000

# Test DNS resolution
docker-compose exec api-gateway nslookup elasticsearch

# Trace network path
docker-compose exec api-gateway traceroute elasticsearch
\end{lstlisting}
\end{commandbox}

\section{Performance Profiling}

\subsection{API Response Time Profiling}

\textbf{Using cURL with Timing:}
\begin{commandbox}
\begin{lstlisting}[style=bashstyle]
# Measure API response time
curl -w "\nTime Total: %{time_total}s\nTime Connect: %{time_connect}s\n" \
  -o /dev/null -s \
  http://localhost:8000/api/v1/documents

# Output:
# Time Total: 0.234s
# Time Connect: 0.012s
\end{lstlisting}
\end{commandbox}

\textbf{Python Profiling:}
\begin{lstlisting}[style=pythonstyle]
# Add profiling to endpoint
from fastapi import Request
import time

@app.middleware("http")
async def add_process_time_header(request: Request, call_next):
    start_time = time.time()
    response = await call_next(request)
    process_time = time.time() - start_time
    response.headers["X-Process-Time"] = str(process_time)
    return response
\end{lstlisting}

\subsection{Database Query Profiling}

\textbf{PostgreSQL Query Analysis:}
\begin{commandbox}
\begin{lstlisting}[style=bashstyle]
# Enable query logging
docker-compose exec postgres psql -U postgres -d document_intelligence -c \
  "ALTER SYSTEM SET log_min_duration_statement = 100;"  # Log queries > 100ms

# Explain query plan
docker-compose exec postgres psql -U postgres -d document_intelligence -c \
  "EXPLAIN ANALYZE SELECT * FROM documents WHERE space_id = 'abc123';"
\end{lstlisting}
\end{commandbox}

\textbf{Elasticsearch Slow Query Log:}
\begin{commandbox}
\begin{lstlisting}[style=bashstyle]
# Enable slow query logging (>1s)
curl -X PUT "http://localhost:9200/documents/_settings" -H 'Content-Type: application/json' -d'
{
  "index.search.slowlog.threshold.query.warn": "1s",
  "index.search.slowlog.threshold.fetch.warn": "1s"
}'

# View slow logs
docker-compose logs elasticsearch | grep "took_millis"
\end{lstlisting}
\end{commandbox}

\section{Getting Help}

\subsection{Support Channels}

\begin{itemize}
    \item \textbf{Technical Documentation:} This manual
    \item \textbf{GitHub Issues:} \texttt{https://github.com/sree181/data-ingestion-platform/issues}
    \item \textbf{Email Support:} \texttt{support@auburn.edu}
    \item \textbf{Community Forum:} (TBD)
\end{itemize}

\subsection{Reporting Bugs}

When reporting bugs, include:
\begin{enumerate}
    \item System environment (OS, Docker version)
    \item Steps to reproduce
    \item Expected vs. actual behavior
    \item Relevant log excerpts
    \item Screenshots (if UI issue)
    \item Error stack traces
\end{enumerate}

\textbf{Example Bug Report:}
\begin{lstlisting}[style=pythonstyle]
Title: Classification fails for documents > 10 pages

Environment:
- OS: Ubuntu 22.04
- Docker: 24.0.6
- Docker Compose: 2.20.0
- RAM: 16 GB

Steps to Reproduce:
1. Upload PDF with 15 pages
2. Enable ICSR classification
3. Monitor processing status

Expected: Document classified within 5 seconds
Actual: Timeout error after 30 seconds

Logs:
[2024-11-01 10:30:15] ERROR: Classification timeout
[2024-11-01 10:30:15] Stack trace: ...

Screenshot: (attached)
\end{lstlisting}

\end{document}

